\documentclass[12pt]{article}
\usepackage[english]{babel}
\usepackage[english]{isodate}
\usepackage[table,svgnames]{xcolor}
\usepackage{url}
\usepackage[utf8x]{inputenc}
\usepackage[T1]{fontenc}
\usepackage{longtable}
\usepackage{amsmath}
\usepackage{graphicx}
\usepackage{parskip}
\usepackage{fancyhdr}
\usepackage{vmargin}
\usepackage{hyperref}
\usepackage{pgfgantt}
\usepackage{pgf-umlcd}
\usepackage{xparse}
\usepackage{float}
\usepackage{tabularx}
\usepackage{titling}
\usepackage[nogroupskip,toc,acronym]{glossaries}
\usepackage{fancyhdr}
\usepackage[authoryear]{natbib}
\usepackage{enumitem}
%\bibliographystyle{fontysIEEtranN}

\pagestyle{fancy}
\fancyhf{}
\lhead{\includegraphics[width=0.2\linewidth]{TreewatchLogo.pdf}}
\rhead{Personal dossier Martijn Bonajo}
\rfoot{Page \thepage}

% alternate row colors for all tables
\definecolor{lightGrey}{rgb}{0.95,0.95,0.95}

\let\oldtable\table
\let\endoldtable\endtable
\renewenvironment{table}{\rowcolors{2}{lightGrey}{}\oldtable}{\endoldtable}

\let\oldtabular\tabular
\let\endoldtabular\endtabular
\renewenvironment{tabular}{\rowcolors{2}{lightGrey}{}\oldtabular}{\endoldtabular}

\let\oldtabularx\tabularx
\let\endoldtabularx\endtabularx
\renewenvironment{tabularx}{\rowcolors{2}{lightGrey}{}\oldtabularx}{\endoldtabularx}

\let\oldlongtable\longtable
\let\endoldlongtable\endlongtable
\renewenvironment{longtable}{\rowcolors{2}{lightGrey}{}\oldlongtable} {\endoldlongtable}

\graphicspath{{../img/}}

\renewcommand{\arraystretch}{1.5}

\DeclareDocumentCommand{\newdualentry}{ O{} O{} m m m m } {
  \newglossaryentry{gls-#3}{name={#5},text={#5\glsadd{#3}},
    description={#6},#1
  }
  \makeglossaries
  \newacronym[see={[Glossary:]{gls-#3}},#2]{#3}{#4}{#5\glsadd{gls-#3}}
}


\newlist{SMART}{description}{2}
\setlist[SMART]{leftmargin=12em,style=nextline}

\newlist{STAR}{description}{2}
\setlist[STAR]{leftmargin=8em,style=nextline}


%\loadglsentries{../glossary.tex}

\setmarginsrb{3 cm}{2.5 cm}{3 cm}{2.5 cm}{1 cm}{1.5 cm}{1 cm}{1.5 cm}

%\makeglossaries
\begin{document}

%%%%%%%%%%%%%%%%%%%%%%%%%%%%%%%%%%%%%%%%%%%%%%%%%%%%%%%%%%%%%%%%%%%%%%%%%%%%%%%%%%%%%%%%% Preface of the report

    \pagenumbering{roman} % Roman numerals for page counter


    % Title Page
    \begin{titlingpage}
        \begin{center}
            \begin{minipage}{\linewidth}
            \centering
            %University logo
            \includegraphics[width=0.3\linewidth]{FontysLogo.pdf}
            \par
            \vspace{3cm}
            %Thesis title
            {\uppercase
                {\Large Personal dossier\\ Martijn Bonajo \\ 2015 Sofa GTL \\ TreeWatch Project
            \par
            \vspace{3cm}}}
            \includegraphics[width=0.5\linewidth]{TreewatchLogo.pdf}
            \par
            \vspace{2cm}
            %Author's name
            {Martijn Bonajo\\
	         \par}
            \vspace{2cm}

            %Date
            \today
            \end{minipage}
        \end{center}
    \end{titlingpage}
    \clearpage

    \section*{Document information}
\addcontentsline{toc}{section}{Document information}
\renewenvironment{tabular}{\oldtabular}{\endoldtabular}
	\begin{tabular}{ll}
		\textbf{Document name:} & Personal dossier Martijn Bonajo\\
		\textbf{Document owner:} & Martijn Bonajo \\
		\textbf{Company/Organisation:} & Fleuren Baarlo \\
		\textbf{Contact person:} & Jan Kerkenhoff \\
		\textbf{Date:} & \today \\
		\textbf{Place:} & Fontys University of Applied Science Venlo \\
		\textbf{Author:} & \parbox[t]{5cm}{
		Martijn Bonajo\\ m.bonajo@student.fontys.nl\\ 2197164 \\}
	\end{tabular}
\renewenvironment{tabular}{\rowcolors{2}{lightGrey}{}\oldtabular}{\endoldtabular}

    %\printglossary[type=\acronymtype]
    %\printglossary
    \pagebreak

    %\listoffigures
    %\addcontentsline{toc}{section}{\listfigurename}
    %\listoftables
    %\addcontentsline{toc}{section}{\listtablename}
    %\pagebreak

    \tableofcontents
    \clearpage

%%%%%%%%%%%%%%%%%%%%%%%%%%%%%%%%%%%%%%%%%%%%%%%%%%%%%%%%%%%%%%%%%%%%%%%%%%%%%%%%%%%%%%%%% Main Part of the Report

    \pagenumbering{arabic}
    
    \section{A-Competencies}
\subsection{SMART Learning Goal 1}
\begin{SMART}
    \item[Specific] As the head developer for this Sofa I want to get more experience developing applications for mobile systems like Android and iOS. Specializing in a framework for running web apps native on different mobile systems
    \item[Measurable] To measure this learning goal I will take into account the number of errors per lines of code.
    \item[Attainable] The topic is widely described in the internet, there are many useful tutorials and guides available, also an extensive documentation is provide for almost all technologies
Relevant
    \item[Relevant] In my future career, in which I want to become a application developer mobile development is often demanded skill.
    \item[Time-limited] By January 2016, this project will be closed, Therefore I will be able to develop mobile for the project until then.
    \item[My complete goal] To improve my mobile development experience by January 2016.
\end{SMART}

\subsection{STAR Analysis 1}
\begin{STAR}
    \item[Situation]
    \item[Task]
    \item[Action]
    \item[Result]
\end{STAR}

\subsection{SMART Learning Goal 2}
\begin{SMART}
    \item[Specific] I wanna improve my skills in presenting my ideas to people.
    \item[Measurable] To measure this learning goal I will take into account the number of different documents, meetings and presentations done for the project.
    \item[Attainable] I already learned a lot in previous projects and the internship. I will try to use the gained knowledge and use that to improve my presentations.
    \item[Relevant] Expressing your thoughts and ideas to other people is always needed in software development.
    \item[Time-limited] By January 2016, this project will be closed, Therefore I will be able improve my communication until then.
    \item[My complete goal] To improve the way I express ideas and thoughts by January 2016.
\end{SMART}

\subsection{STAR Analysis 2}
\begin{STAR}
    \item[Situation]
    \item[Task]
    \item[Action]
    \item[Result]
\end{STAR}

\section{B-Competencies}

\subsection{SMART Learning Goal 3}
\begin{SMART}
    \item[Specific] I wanna improve my skills in presenting my ideas to people.
    \item[Measurable] To measure this learning goal I will take into account the number of different documents, meetings and presentations done for the project.
    \item[Attainable] I already learned a lot in previous projects and the internship. I will try to use the gained knowledge and use that to improve my presentations.
    \item[Relevant] Expressing your thoughts and ideas to other people is always needed in software development.
    \item[Time-limited] By January 2016, this project will be closed, Therefore I will be able improve my communication until then.
    \item[My complete goal] To improve the way I express ideas and thoughts by January 2016.
\end{SMART}

\subsection{STAR Analysis 3}
\begin{STAR}
    \item[Situation]
    \item[Task]
    \item[Action]
    \item[Result]
\end{STAR}


%\clearpage

%\bibliography{biblist}

\end{document}
